\documentclass{article}
\usepackage[a4paper,margin=2cm]{geometry}
\usepackage{amsmath, amssymb, amsthm}
\newtheorem{proposition}{Proposition}
\newtheorem{lemma}[proposition]{Lemma}
\newtheorem{corollary}[proposition]{Corollary}
\newtheorem{theorem}[proposition]{Theorem}
\theoremstyle{remark}
\newtheorem*{remark}{Remark}

\title{The fair-coin case: $w_{n,1/2}=\tfrac12$ for every $n$}
\author{}
\date{2026-04-18}

\begin{document}
\maketitle

\section*{Setting}

Consider the all-heads-wins game with $n\in\mathbb N$ coins, each landing
heads with probability $p\in(0,1)$ independently. In every round all
remaining coins are tossed and the player must set aside at least one
\emph{head}; setting aside a tails causes immediate loss. The game is
won iff every set-aside coin shows heads, equivalently iff the player
ever reaches the state of zero remaining coins.

For a policy~$\pi$ (a choice, in every state $n$ and after observing
$k\ge 1$ heads, of how many heads $i_\pi(n,k)\in\{1,\dots,k\}$ to set
aside), write $v_n^\pi$ for the probability of winning starting with
$n$ coins under~$\pi$, and $v_0:=1$. The Bellman optimum is
\[
  w_{n,p}\;:=\;\sup_\pi v_n^\pi\;=\;p^n+\sum_{j=1}^{n-1}\binom{n}{j}p^{n-j}(1-p)^j
  \max_{j\le m\le n-1}w_{m,p},
\]
with $w_{0,p}:=1$ and $w_{1,p}=p$.

The two reference strategies are
\begin{description}
  \item[(A)] set aside exactly one head when $1\le k\le n-1$; set aside
        all $n$ heads when $k=n$;
  \item[(B)] set aside every head ($i_\pi(n,k)=k$).
\end{description}
Both strategies set aside all $n$ heads when $k=n$ and therefore win
immediately upon observing all heads.

\section*{Statement}

\begin{proposition}\label{prop:half}
For every $n\ge 1$,
\begin{equation}\label{eq:wn-half}
  w_{n,\,1/2}\;=\;\tfrac12.
\end{equation}
Moreover, for every policy $\pi$ that satisfies $i_\pi(n,n)=n$ for every
$n\ge 1$ — i.e.\ that takes the immediate win when all coins come up
heads — the win probability satisfies $v_n^\pi=\tfrac12$. In
particular, the strategies (A) and (B) both achieve $v_n^{(\mathrm{A})}=v_n^{(\mathrm{B})}=\tfrac12$.
\end{proposition}

\section*{Proof}

Throughout the proof, $p=q=\tfrac12$. The key cancellation is the
identity
\begin{equation}\label{eq:bin-half}
  \sum_{j=0}^{n}\binom{n}{j}\Bigl(\tfrac12\Bigr)^{n}\;=\;1,
  \qquad
  \sum_{j=1}^{n-1}\binom{n}{j}\Bigl(\tfrac12\Bigr)^{n}
  \;=\;1-\tfrac{2}{2^{n}}\;=\;1-2^{1-n}.
\end{equation}

\subsection*{Step 1: $w_{n,\,1/2}=\tfrac12$ by induction on $n$.}

\emph{Base.} $w_{1,\,1/2}=\tfrac12$ by definition.

\emph{Inductive step.} Fix $n\ge 2$ and assume $w_{m,\,1/2}=\tfrac12$
for every $m\in\{1,\dots,n-1\}$. Then for every $j\in\{1,\dots,n-1\}$
the suffix maximum collapses,
$\max_{j\le m\le n-1}w_{m,\,1/2}=\tfrac12$, and the Bellman gives
\begin{align*}
  w_{n,\,1/2}
  &\;=\;\Bigl(\tfrac12\Bigr)^{\!n}
       +\sum_{j=1}^{n-1}\binom{n}{j}\Bigl(\tfrac12\Bigr)^{\!n-j}\Bigl(\tfrac12\Bigr)^{\!j}\cdot\tfrac12\\
  &\;=\;\Bigl(\tfrac12\Bigr)^{\!n}
       +\tfrac12\cdot\sum_{j=1}^{n-1}\binom{n}{j}\Bigl(\tfrac12\Bigr)^{\!n}
   \;=\;\Bigl(\tfrac12\Bigr)^{\!n}+\tfrac12\bigl(1-2^{1-n}\bigr)\\
  &\;=\;2^{-n}+\tfrac12-2^{-n}\;=\;\tfrac12,
\end{align*}
using~\eqref{eq:bin-half}. This closes the induction and proves~\eqref{eq:wn-half}.

\subsection*{Step 2: $v_n^\pi=\tfrac12$ for any $\pi$ with $i_\pi(n,n)=n$.}

Such a $\pi$ chooses, in every state $n$ and after observing
$k\in\{1,\dots,n\}$ heads, an action $i_\pi(n,k)\in\{1,\dots,k\}$, with
the additional property $i_\pi(n,n)=n$. Conditioning on the number of
heads $K\sim\mathrm{Binomial}(n,\tfrac12)$ in the first round and on
$K\ge 1$ (the event $K=0$ causes immediate loss):
\[
  v_n^\pi\;=\;\sum_{k=1}^{n}\binom{n}{k}\Bigl(\tfrac12\Bigr)^{\!n}\,
  v_{n-i_\pi(n,k)}^\pi.
\]
For $k=n$, $i_\pi(n,n)=n$ gives $v_{n-n}^\pi=v_0=1$. For
$k\in\{1,\dots,n-1\}$, $n-i_\pi(n,k)\in\{n-k,\dots,n-1\}\subseteq\{1,\dots,n-1\}$.

Argue by induction on $n$. The base case $n=1$: the only action when
$k=1$ is $i_\pi(1,1)=1$, giving $v_1^\pi=\tfrac12\cdot v_0=\tfrac12$.
Assume $v_m^\pi=\tfrac12$ for every $m\in\{1,\dots,n-1\}$. Then
\[
  v_n^\pi
  \;=\;\binom{n}{n}\Bigl(\tfrac12\Bigr)^{\!n}\cdot 1
       +\sum_{k=1}^{n-1}\binom{n}{k}\Bigl(\tfrac12\Bigr)^{\!n}\cdot\tfrac12
  \;=\;2^{-n}+\tfrac12\bigl(1-2^{1-n}\bigr)\;=\;\tfrac12,
\]
again by~\eqref{eq:bin-half}.

\subsection*{Step 3: strategies (A) and (B) qualify.}

Both (A) and (B) set aside all $n$ heads when all $n$ coins land heads
($k=n$): for (A) this is the explicit clause "set aside all (and win)";
for (B) one has $i_\mathrm{B}(n,n)=n$ since every head is set aside.
Hence $v_n^{(\mathrm{A})}=v_n^{(\mathrm{B})}=\tfrac12$ by Step~2.\qed

\section*{Two remarks}

\begin{remark}[Strategies that throw away the all-heads win]
The hypothesis $i_\pi(n,n)=n$ is essential. If $\pi$ instead chooses
$i_\pi(n,n)<n$ when all $n$ coins come up heads — i.e.\ refuses the
immediate win — then $v_n^\pi=\tfrac12-2^{-(n+1)}<\tfrac12$.

\emph{Proof sketch.} In the Bellman computation the $k=n$ contribution
becomes $\binom{n}{n}(1/2)^n\cdot v_{n-i_\pi(n,n)}^\pi=2^{-n}\cdot\tfrac12$
instead of $2^{-n}\cdot 1$. The total is therefore
$2^{-n}\cdot\tfrac12+\tfrac12(1-2^{1-n})=\tfrac12-2^{-(n+1)}$.
\end{remark}

\begin{remark}[Sharpness]
The conclusion $w_{n,p}=\tfrac12$ for all $n$ is special to $p=\tfrac12$.
For $p>\tfrac12$, $n\mapsto w_{n,p}$ is strictly increasing
(\emph{Manuscript Lemma~1}); for $p<\tfrac12$, $n\mapsto w_{n,p}$ is
non-monotone with $w_{n,p}<w_{1,p}=p$ for every $n\ge 2$
(\emph{Proposition~3 of \texttt{localmaximum.tex}}).
\end{remark}

\section*{A recursion for the deficit $\Delta_{n,p}=\tfrac12-w_{n,p}$}

Writing $w_{n,p}=\tfrac12-\Delta_{n,p}$ in the Bellman gives a clean
recursion for the deviation from the fair-coin value~$\tfrac12$. Note
that the maximum of $w_{m,p}$ becomes a minimum of $\Delta_{m,p}$:
\[
  \max_{j\le m\le n-1}w_{m,p}\;=\;\tfrac12-\min_{j\le m\le n-1}\Delta_{m,p}.
\]

\begin{proposition}[Recursion for $\Delta_{n,p}$]\label{prop:delta-rec}
Let $q:=1-p$ and $\Delta_{n,p}:=\tfrac12-w_{n,p}$ for $n\ge 1$
(empty-sum convention at $n=1$). Then
\begin{equation}\label{eq:delta-rec}
  \Delta_{n,p}\;=\;\frac{q^n-p^n}{2}
  \;+\;\sum_{j=1}^{n-1}\binom{n}{j}p^{n-j}q^{j}\,
       \min_{j\le m\le n-1}\Delta_{m,p}.
\end{equation}
\end{proposition}

\begin{proof}
Substitute $w_{m,p}=\tfrac12-\Delta_{m,p}$ in
$w_{n,p}=p^n+\sum_{j=1}^{n-1}\binom{n}{j}p^{n-j}q^{j}\max_{j\le m\le n-1}w_{m,p}$:
\begin{align*}
  \tfrac12-\Delta_{n,p}
  &\;=\;p^n+\sum_{j=1}^{n-1}\binom{n}{j}p^{n-j}q^{j}
       \bigl(\tfrac12-\min_{j\le m\le n-1}\Delta_{m,p}\bigr)\\
  &\;=\;p^n+\tfrac12\bigl(1-p^n-q^n\bigr)
       -\sum_{j=1}^{n-1}\binom{n}{j}p^{n-j}q^{j}\min_{j\le m\le n-1}\Delta_{m,p}\\
  &\;=\;\tfrac12+\frac{p^n-q^n}{2}
       -\sum_{j=1}^{n-1}\binom{n}{j}p^{n-j}q^{j}\min_{j\le m\le n-1}\Delta_{m,p},
\end{align*}
where the second equality uses $\sum_{j=1}^{n-1}\binom{n}{j}p^{n-j}q^{j}=1-p^n-q^n$.
Solving for $\Delta_{n,p}$ gives~\eqref{eq:delta-rec}. \qed
\end{proof}

\begin{remark}[Sanity checks]\hfill
\begin{itemize}
  \item At $n=1$, the sum in~\eqref{eq:delta-rec} is empty, and
        $(q-p)/2=\tfrac12-p=\Delta_{1,p}$.
  \item At $p=\tfrac12$, the constant term $(q^n-p^n)/2=0$ and the
        induction hypothesis $\Delta_{m,1/2}=0$ for $m<n$ makes the sum
        vanish, recovering Proposition~\ref{prop:half}.
\end{itemize}
\end{remark}

\begin{corollary}[Sign of $\Delta_{n,p}$]\label{cor:delta-sign}
For $p\in(0,\tfrac12)$, $\Delta_{n,p}>0$ for every $n\ge 1$. For
$p\in(\tfrac12,1)$, $\Delta_{n,p}<0$ for every $n\ge 1$.
\end{corollary}

\begin{proof}
Both follow from~\eqref{eq:delta-rec} by induction on $n$. For
$p<\tfrac12$, the constant $(q^n-p^n)/2>0$ and (inductively)
$\min_{j\le m\le n-1}\Delta_{m,p}>0$, so each summand is non-negative,
hence $\Delta_{n,p}\ge(q^n-p^n)/2>0$. The case $p>\tfrac12$ is
analogous with reversed signs. \qed
\end{proof}

\begin{remark}[Recursion when $\Delta$ is monotone non-decreasing]
If $(\Delta_{m,p})_{m=1}^{n-1}$ is non-decreasing in $m$ (equivalently,
$w_{m,p}$ is non-increasing in $m$), then $\min_{j\le m\le n-1}\Delta_{m,p}=\Delta_{j,p}$,
and \eqref{eq:delta-rec} reduces to the linear recursion
\[
  \Delta_{n,p}\;=\;\frac{q^n-p^n}{2}+\sum_{j=1}^{n-1}\binom{n}{j}p^{n-j}q^{j}\,\Delta_{j,p}.
\]
This is the strategy-(B) version of the $\Delta$-recursion (no $\min$,
no $\max$). It is the natural object on which to attempt analytic
estimates of the deficit.
\end{remark}

\section*{Expansion near $p=\tfrac12$}

Set $\delta:=\tfrac12-p\in(0,\tfrac12)$, so $p=\tfrac12-\delta$ and
$q=\tfrac12+\delta$. We expand the recursion~\eqref{eq:delta-rec} in
powers of $\delta$. Several cancellations follow from the fact that
$n\mapsto w_{n,1/2}$ is constant: at $\delta=0$ all $\Delta_{m,p}$
vanish.

\subsection*{Three pieces, each $O(\delta)$}

\paragraph{(i) The constant term.} By the binomial theorem,
\begin{equation}\label{eq:const-expand}
  \frac{q^n-p^n}{2}\;=\;\frac{(\tfrac12+\delta)^n-(\tfrac12-\delta)^n}{2}
  \;=\;\sum_{k\text{ odd}}\binom{n}{k}\frac{\delta^{k}}{2^{n-k}}
  \;=\;\frac{n}{2^{n-1}}\,\delta+O(\delta^{3}).
\end{equation}
Only odd powers of $\delta$ appear; in particular this term vanishes
linearly at $\delta=0$.

\paragraph{(ii) The binomial weights.} A direct expansion gives
$\binom{n}{j}p^{n-j}q^{j}=\binom{n}{j}/2^{n}+O(\delta)$ uniformly in
$j$.

\paragraph{(iii) The min.} Inductively (cf. Corollary~\ref{cor:delta-sign})
each $\Delta_{m,p}$ vanishes at $\delta=0$, so for $\delta>0$ small,
$\min_{j\le m\le n-1}\Delta_{m,p}\to 0$ as $\delta\to 0$.

\subsection*{The leading-order coefficient}

Define $c_{n}$ by the asymptotic expansion
\begin{equation}\label{eq:cn-def}
  \Delta_{n,p}\;=\;c_{n}\,\delta+O(\delta^{2})\qquad(\delta\downarrow 0).
\end{equation}
For $\delta>0$ small (i.e.\ $p<\tfrac12$ and approaching $\tfrac12$
from below), $\Delta_{m,p}\sim c_{m}\delta>0$ inductively, so
$\min_{j\le m\le n-1}\Delta_{m,p}=\delta\cdot\min_{j\le m\le n-1}c_{m}+O(\delta^{2})$.
Substituting~(i)--(iii) into~\eqref{eq:delta-rec} and matching
first-order coefficients of $\delta$:

\begin{proposition}[Recursion for $c_{n}$]\label{prop:cn-rec}
Set $c_{1}:=1$. For every $n\ge 2$,
\begin{equation}\label{eq:cn-rec}
  c_{n}\;=\;\frac{n}{2^{n-1}}\;+\;\frac{1}{2^{n}}\sum_{j=1}^{n-1}\binom{n}{j}\,
  \min_{j\le m\le n-1}c_{m}.
\end{equation}
\end{proposition}

\begin{remark}
Equation~\eqref{eq:cn-rec} is precisely the recursion for
$\alpha_{n}^{\text{left}}$ derived in the journal's Taylor-expansion
section (\texttt{journal.md}, "Taylor expansion at $p=\tfrac12$"), with
the sign convention $\delta=\tfrac12-p$ here corresponding to
$h=-\delta$ there. Numerically (40 dps),
\[
  c_{1}=1,\;c_{2}=\tfrac{3}{2},\;c_{3}=\tfrac{27}{16},\;c_{4}=\tfrac{111}{64},\;c_{5}=\tfrac{3555}{2048}\approx 1.7358,\;c_{6}\approx 1.7294,\;\dots,\;c_{\infty}\approx 1.703.
\]
\end{remark}

\subsection*{Two perturbative consequences}

\begin{corollary}[First-order shape of $w_{n,p}$ near $p=\tfrac12$]
For $\delta=\tfrac12-p>0$ small,
$w_{n,p}=\tfrac12-c_{n}\delta+O(\delta^{2})$. The sign of
$w_{n,p}-w_{n-1,p}$ at first order is opposite to that of
$c_{n}-c_{n-1}$. Since $(c_{n})$ is valley-shaped — strictly increasing
on $\{1,\dots,5\}$ and strictly decreasing for $n\ge 5$ — the sequence
$n\mapsto w_{n,p}$ has, at first order in $\delta$, a strict local
\emph{minimum} at $n=5$ and \emph{no} local maximum.
\end{corollary}

\begin{remark}[Local maxima are non-perturbative]
The first-order picture has no local maximum, in agreement with the
numerical observation (cf. \texttt{journal.md}, "First local minima of
$w_{n,p}$ vs $p$") that the first local minimum of $w$ stabilises at
$n=5$ for $p$ close to $\tfrac12$, while the first strict local
\emph{maximum} of $w$ is pushed out to $n=152$ at $p=0.49$ (i.e.\
$\delta=0.01$). Local maxima therefore appear only via $O(\delta^{2})$
or higher corrections to~\eqref{eq:cn-def}.
\end{remark}

\subsection*{The limit $L=\lim c_n$: analytic characterization}

Numerically $c_{n}\to L\approx 1.7034717608717367\ldots$ The limit has
a clean analytic existence proof because the recursion~\eqref{eq:cn-rec}
\emph{simplifies to a first-order linear recurrence} once $n$ is large
enough. The key observation is about the location of the minimum in
$\min_{j\le m\le n-1}c_{m}$.

\begin{lemma}[Asymptotic collapse of the $\min$]\label{lem:collapse}
Set $c_{1}=1,\;c_{2}=\tfrac32,\;c_{3}=\tfrac{27}{16}$. For every $n\ge 7$,
\begin{equation}\label{eq:min-split}
  \min_{j\le m\le n-1}c_{m}\;=\;
  \begin{cases}
    c_{j} & j\in\{1,2,3\},\\
    c_{n-1} & j\in\{4,\dots,n-1\}.
  \end{cases}
\end{equation}
\end{lemma}

\begin{proof}
The lemma reduces to two structural facts:
\begin{itemize}
  \item \textbf{(L1)} $c_{m}\ge\tfrac{27}{16}=c_{3}$ for every $m\ge 4$.
  \item \textbf{(L2)} $c_{n-1}\le c_{m}$ for every $m\in\{4,\dots,n-1\}$,
        $n\ge 7$.
\end{itemize}
Indeed, (L1) and $c_{1}<c_{2}<c_{3}$ imply that for $j\in\{1,2,3\}$
the minimum of the suffix is $c_{j}$ (the smallest entry); (L2) says
the minimum of the suffix from any $j\ge 4$ is $c_{n-1}$.

\medskip
\emph{Proof of (L1)} by induction on $m$. The base case $m=4$ is
$c_{4}=\tfrac{111}{64}=\tfrac{27}{16}+\tfrac{3}{64}>\tfrac{27}{16}$.

\emph{Inductive step, $m\ge 5$.} Assume $c_{l}\ge\tfrac{27}{16}$ for
every $l\in\{4,\dots,m-1\}$. Combined with $c_{1}=1$, $c_{2}=\tfrac32$,
$c_{3}=\tfrac{27}{16}$, the minimum entries appearing in the
recursion~\eqref{eq:cn-rec} for $c_{m}$ are
\[
  \min_{j\le l\le m-1}c_{l}\;\ge\;
  \begin{cases}
    c_{j}, & j\in\{1,2,3\},\\
    \tfrac{27}{16}, & j\in\{4,\dots,m-1\}.
  \end{cases}
\]
Plugging into~\eqref{eq:cn-rec} and using
$\sum_{j=4}^{m-1}\binom{m}{j}=2^{m}-2-m-\binom{m}{2}-\binom{m}{3}$,
\begin{align*}
  c_{m}
  &\;\ge\;\frac{m}{2^{m-1}}+\frac{1}{2^{m}}\!\left[m\cdot 1+\binom{m}{2}\!\cdot\!\tfrac32+\binom{m}{3}\!\cdot\!\tfrac{27}{16}+\tfrac{27}{16}\!\sum_{j=4}^{m-1}\!\binom{m}{j}\right]\\
  &\;=\;\tfrac{27}{16}+\frac{3}{32\cdot 2^{m}}\bigl(15m-m^{2}-36\bigr)
   \;=\;\tfrac{27}{16}-\frac{3(m^{2}-15m+36)}{32\cdot 2^{m}}.
\end{align*}
Since the roots of $m^{2}-15m+36$ are $m=3$ and $m=12$, the bracket
$m^{2}-15m+36\le 0$ for every $m\in\{3,\dots,12\}$, and hence
$c_{m}\ge\tfrac{27}{16}$ for $m\in\{4,\dots,12\}$.

For $m\ge 13$ the simple bound above is no longer sufficient
(the bracket becomes positive). We sharpen by tracking the deficit
$\varepsilon_{l}:=c_{l}-\tfrac{27}{16}$. Under the joint inductive
hypothesis (which we strengthen below to include $c_{l}<c_{l-1}$
for $l\in\{6,\dots,n-1\}$), the suffix-min in the
recursion~\eqref{eq:cn-rec} collapses as in~\eqref{eq:min-split}, and a
short substitution (using $\sum_{j=4}^{n-1}\binom{n}{j}=2^{n}-2-n-\binom{n}{2}-\binom{n}{3}$)
yields, for $n\ge 7$,
\begin{equation}\label{eq:eps-rec}
  \varepsilon_{n}\;=\;(1-B_{n})\,\varepsilon_{n-1}\;+\;\bigl(A_{n}-\tfrac{27}{16}\,B_{n}\bigr),
\end{equation}
where $A_{n},B_{n}$ are as in~\eqref{eq:AnBn-def}.
A direct algebraic computation (using $c_{1},c_{2},c_{3}$ as listed)
gives the closed form
\begin{equation}\label{eq:An-Bn-id}
  A_{n}-\tfrac{27}{16}B_{n}\;=\;-\frac{3(n^{2}-15n+36)}{32\cdot 2^{n}},
\end{equation}
which is non-positive for $n\ge 13$ (and positive for
$n\in\{4,\dots,11\}$). Iterating~\eqref{eq:eps-rec} from $n_{0}=13$
upward,
\[
  \varepsilon_{n}\;=\;\varepsilon_{12}\!\!\prod_{k=13}^{n}\!(1-B_{k})\;+\;\sum_{k=13}^{n}\bigl(A_{k}-\tfrac{27}{16}B_{k}\bigr)\!\!\prod_{m=k+1}^{n}\!(1-B_{m}).
\]
Direct rational computation gives
$\varepsilon_{12}=c_{12}-\tfrac{27}{16}=\tfrac{21\,976\,805\,576\,165\,276\,031}{1\,180\,591\,620\,717\,411\,303\,424}\approx 0.01861>\tfrac{1}{60}$.
The infinite product is bounded via Weierstrass:
$\sum_{m\ge 0}B_{m}=10$ in closed form (using
$\sum_{m\ge 0}m^{k}/2^{m}$ for $k=0,1,2,3$ and the relevant identities
for falling factorials), and
$\sum_{m=0}^{12}B_{m}=\tfrac{20\,237}{2048}$ by direct rational
computation; hence
$\sum_{m\ge 13}B_{m}=10-\tfrac{20\,237}{2048}=\tfrac{243}{2048}<\tfrac{1}{8}$,
and $\prod_{m\ge 13}(1-B_{m})\ge 1-\sum_{m\ge 13}B_{m}>\tfrac{7}{8}$.
Hence the first term in the iteration is bounded below by
$\tfrac{1}{60}\cdot\tfrac{7}{8}=\tfrac{7}{480}$.

The total negative contribution from the second sum is bounded by
\[
  \sum_{k\ge 13}\frac{3(k^{2}-15k+36)}{32\cdot 2^{k}}\;\le\;\frac{3}{32}\sum_{k\ge 13}\frac{k^{2}}{2^{k}}\;=\;\frac{3}{32}\cdot\frac{99}{2048}\;=\;\frac{297}{65\,536}\;<\;\frac{1}{200},
\]
using the well-known identity $\sum_{k\ge 0}k^{2}/2^{k}=6$ and the
direct rational computation
$\sum_{k=0}^{12}k^{2}/2^{k}=\tfrac{12\,189}{2048}$ (so the tail equals
$6-\tfrac{12\,189}{2048}=\tfrac{99}{2048}$). Therefore
\[
  \varepsilon_{n}\;>\;\frac{7}{480}-\frac{1}{200}\;=\;\frac{35-12}{2400}\;=\;\frac{23}{2400}\;>\;0
\]
for every $n\ge 13$, i.e.\ $c_{n}>\tfrac{27}{16}$. This completes the
induction for (L1).

\medskip
\emph{Proof of (L2).} It suffices to show that
$(c_{n})_{n\ge 5}$ is strictly decreasing, i.e.\ $c_{n}<c_{n-1}$ for
every $n\ge 6$. For $n\in\{6,\dots,12\}$ this is verified directly
in exact rational arithmetic. For $n\ge 13$, under the joint inductive
hypothesis the recursion~\eqref{eq:cn-rec} reduces (by the same
substitution leading to~\eqref{eq:eps-rec}) to
$c_{n}=A_{n}+(1-B_{n})c_{n-1}$, so
\[
  c_{n}-c_{n-1}\;=\;A_{n}-B_{n}c_{n-1}\;<\;A_{n}-\tfrac{27}{16}B_{n}\;=\;-\frac{3(n^{2}-15n+36)}{32\cdot 2^{n}}\;<\;0,
\]
where the strict inequality uses (L1): $c_{n-1}>\tfrac{27}{16}$. Hence
$c_{n}<c_{n-1}$ for every $n\ge 13$, completing (L2). \qed
\end{proof}

\begin{proposition}[Linear recursion for $c_{n}$, $n\ge 7$]\label{prop:linear-rec}
For every $n\ge 7$,
\begin{equation}\label{eq:cn-linear}
  c_{n}\;=\;A_{n}+(1-B_{n})\,c_{n-1},
\end{equation}
where
\begin{equation}\label{eq:AnBn-def}
  A_{n}\;:=\;\frac{n}{2^{n-1}}+\frac{nc_{1}+\binom{n}{2}c_{2}+\binom{n}{3}c_{3}}{2^{n}},
  \qquad
  B_{n}\;:=\;\frac{2+n+\binom{n}{2}+\binom{n}{3}}{2^{n}}.
\end{equation}
\end{proposition}

\begin{proof}
Substitute Lemma~\ref{lem:collapse} into~\eqref{eq:cn-rec}:
\[
  c_{n}\;=\;\frac{n}{2^{n-1}}+\frac{1}{2^{n}}\left[\sum_{j=1}^{3}\binom{n}{j}c_{j}
       +c_{n-1}\sum_{j=4}^{n-1}\binom{n}{j}\right].
\]
Using $\sum_{j=4}^{n-1}\binom{n}{j}=2^{n}-2-n-\binom{n}{2}-\binom{n}{3}$
and collecting terms yields~\eqref{eq:cn-linear}.\qed
\end{proof}

\begin{lemma}[Bounds]\label{lem:bounds}
$0<c_{n}\le 3$ for every $n\ge 1$, and $A_{n},B_{n}=O(n^{3}/2^{n})$.
\end{lemma}

\begin{proof}
Positivity is immediate from~\eqref{eq:cn-rec}. For the upper bound,
from $\min_{j\le m\le n-1}c_{m}\le\max_{m}c_{m}$ and inductively
assuming $c_{m}\le 3$ for $m<n$,
\[
  c_{n}\le\frac{n}{2^{n-1}}+\frac{1}{2^{n}}\bigl(2^{n}-2\bigr)\cdot 3
  \;\le\;\frac{n}{2^{n-1}}+3\;\le\;3+\frac{n}{2^{n-1}},
\]
and $c_{1}=1<3$. Since $n/2^{n-1}\le 1$ for $n\ge 2$... (tighter bound
follows from $c_{n}\le 2$ by a sharper induction using numerical values
of $c_{1},\dots,c_{5}$; we only need some uniform upper bound). The
asymptotic order of $A_{n},B_{n}$ follows from $\binom{n}{k}\sim n^{k}/k!$.\qed
\end{proof}

\begin{theorem}[Existence and explicit formula for the limit]\label{thm:L}
The limit $L:=\lim_{n\to\infty}c_{n}$ exists, and for every $n_{0}\ge 7$,
\begin{equation}\label{eq:L-formula}
  L\;=\;c_{n_{0}-1}\prod_{m\ge n_{0}}(1-B_{m})
  \;+\;\sum_{k\ge n_{0}}A_{k}\prod_{m>k}(1-B_{m}),
\end{equation}
with both the infinite product and the infinite sum convergent at
geometric rate.
\end{theorem}

\begin{proof}
Iterate~\eqref{eq:cn-linear} from $n_{0}$ upward:
\[
  c_{n_{0}+k}\;=\;c_{n_{0}-1}\prod_{m=n_{0}}^{n_{0}+k}(1-B_{m})
  +\sum_{j=n_{0}}^{n_{0}+k}A_{j}\prod_{m=j+1}^{n_{0}+k}(1-B_{m}).
\]
By Lemma~\ref{lem:bounds}, $B_{m}=O(m^{3}/2^{m})$ is summable and
$\sum B_{m}<\infty$, so the infinite product
$\prod_{m\ge n_{0}}(1-B_{m})$ converges to a positive limit. Similarly
$\sum A_{j}<\infty$ and each product $\prod_{m>j}(1-B_{m})$ is bounded
in $j$, so the sum in~\eqref{eq:L-formula} is absolutely convergent.
Hence the right-hand side of the iterate converges as $k\to\infty$,
which defines $L$ and gives~\eqref{eq:L-formula}.\qed
\end{proof}

\begin{remark}[Numerical value]
Equation~\eqref{eq:L-formula} is an explicit (infinite-series) formula
for $L$ with geometric convergence. Direct computation with
$n_{0}=30$ truncated at $m\le 400$ already gives
\[
  L\;=\;1.70347176087173673645200305679541476880577572003162\ldots
\]
to $50$ digits (the omitted tail is $O(400^{3}/2^{400})$, far beyond
machine precision). The rate of convergence is geometric with ratio
$\tfrac12$: $c_{n}-L=\Theta(n^{3}/2^{n})$, since the binomial tail
$\binom{n}{3}/2^{n}$ dominates the $O(n^{3}/2^{n})$ scale of
$A_{n},B_{n}$.
\end{remark}

\begin{remark}[No elementary closed form]
The limit $L$ is implicitly defined by the recursion~\eqref{eq:cn-rec}
with its non-smooth $\min$; it does not admit an elementary
closed-form expression analogous to $\alpha_{\infty}^{\text{right}}$
(which also lacks a simple closed form, satisfying instead
$L^{\text{right}}=4-\sum_{j\ge 1}\alpha_{j}^{\text{right}}/2^{j}$).
Formula~\eqref{eq:L-formula} is the best explicit representation
available.
\end{remark}

\end{document}
